% =====================================================================
% Mục 3. Tiền xử lý số liệu
%
% Gói cần có trong file chính: booktabs, tabularx, array, listings, xcolor,
%   tikz (thư viện positioning, arrows.meta), caption.
% Trước khi \input file này, file chính phải nạp các con số do code sinh ra:
%   % Generated by scripts/report_preprocessing.py; do not edit by hand.
\newcommand{\PreRowsRaw}{3406}
\newcommand{\PreColsRaw}{34}
\newcommand{\PreColsSelected}{18}
\newcommand{\PreCoreSpeedDash}{907}
\newcommand{\PreCoreSpeedBlank}{29}
\newcommand{\PreUnknownReleaseDate}{30}
\newcommand{\PreBandwidthInMb}{4}
\newcommand{\PreRopsMultiplier}{305}
\newcommand{\PreTextureMismatchBefore}{274}
\newcommand{\PreTextureMismatchAfter}{0}
\newcommand{\PreIntegrated}{374}
\newcommand{\PreUnknownDedicated}{14}
\newcommand{\PreRowsDedicated}{3018}
\newcommand{\PreRowsSingleGpu}{2775}
\newcommand{\PreRowsReleased}{2762}
\newcommand{\PreRowsFinal}{2762}
\newcommand{\PreColsFinal}{16}
\newcommand{\PreDesktop}{2318}
\newcommand{\PreLaptop}{444}
\newcommand{\PrePowerMissing}{481}
\newcommand{\PrePowerMissingPct}{17,4}
\newcommand{\PreMissingMedianYear}{2011}
\newcommand{\PreObservedMedianYear}{2013}
\newcommand{\PreMissingYearP}{$<$ 0,001}
\newcommand{\PreMissingAmdPct}{53,1}
\newcommand{\PreObservedAmdPct}{39,0}
\newcommand{\PreMissingMakerP}{$<$ 0,001}
\newcommand{\PreMissingPctGpuId}{0,0}
\newcommand{\PreMissingPctName}{0,0}
\newcommand{\PreMissingPctManufacturer}{0,0}
\newcommand{\PreMissingPctSegment}{0,0}
\newcommand{\PreMissingPctMemoryType}{0,7}
\newcommand{\PreMissingPctMemoryTypeGroup}{0,7}
\newcommand{\PreMissingPctArchitecture}{1,5}
\newcommand{\PreMissingPctReleaseYear}{1,1}
\newcommand{\PreMissingPctMaxPowerW}{17,4}
\newcommand{\PreMissingPctNeedsAuxPower}{17,4}
\newcommand{\PreMissingPctTextureRatePerWatt}{24,2}
\newcommand{\PreMissingPctTmus}{13,0}
\newcommand{\PreMissingPctCoreSpeedMhz}{22,0}
\newcommand{\PreMissingPctMemoryBandwidthGbS}{1,3}
\newcommand{\PreMissingPctMemoryMb}{0,9}
\newcommand{\PreMissingPctProcessNm}{10,5}

% (tạo bằng: python scripts/report_preprocessing.py; script này cũng sinh
% generated/preprocessing_sample_rows.tex cho bảng dữ liệu mẫu). Mọi con số
% trong lời văn đều là macro \Pre..., nên chạy lại code là báo cáo tự cập nhật.
% Không tham chiếu số mục hay số bảng của mục khác (cấu trúc báo cáo chưa chốt);
% gọi theo nội dung: "phần tổng quan", "phần hồi quy", "các mục sau".
% =====================================================================

\section{Tiền xử lý số liệu}
\label{sec:preprocessing}

Phần này trình bày quá trình chuyển tệp \texttt{All\_GPUs.csv} (\PreRowsRaw{} dòng,
\PreColsRaw{} biến, mô tả ở phần tổng quan) thành bộ dữ liệu dùng cho
các phân tích ở mục sau. Toàn bộ quy trình viết bằng Python (thư viện
\texttt{pandas}), chạy lại được từ đầu và cho cùng một kết quả. Quy trình gồm sáu bước:
đọc dữ liệu, chọn biến, làm sạch chuỗi, xử lý định dạng, lọc quan sát và tạo biến mới
(Hình~\ref{fig:preprocess}).

Nhóm tuân theo ba nguyên tắc: không sửa tay dữ liệu gốc; mỗi hàm đọc giá trị chỉ chấp
nhận đúng định dạng đã biết và dừng chương trình nếu gặp giá trị lạ; không điền giá trị
cho biến mục tiêu \texttt{Max\_Power}. Ngoại lai và biến đổi logarit được xử lý trong
từng phân tích ở các mục sau, vì mỗi phân tích cần một cách xử lý khác nhau.

\begin{figure}[!htb]
  \centering
  \begin{tikzpicture}[
      node distance=3.5mm,
      box/.style={rectangle, rounded corners, draw, thick, fill=blue!6,
        text width=7.5cm, align=center, minimum height=7mm, font=\footnotesize},
      io/.style={rectangle, draw, thick, fill=orange!12, text width=7.5cm,
        align=center, minimum height=7mm, font=\footnotesize},
      arrow/.style={-Stealth, thick}]
    \node[io] (raw) {Tệp gốc \texttt{All\_GPUs.csv}: \PreRowsRaw{} dòng $\times$ \PreColsRaw{} biến};
    \node[box, below=of raw] (read) {1. Đọc mọi cột dưới dạng chuỗi};
    \node[box, below=of read] (select) {2. Chọn \PreColsSelected{} biến cần cho đề tài};
    \node[box, below=of select] (clean) {3. Xóa khoảng trắng, đánh dấu giá trị thiếu};
    \node[box, below=of clean] (parse) {4. Tách đơn vị, đổi kiểu, sửa lỗi ``(x$N$)''};
    \node[box, below=of parse] (filter) {5. Lọc: GPU rời, một GPU, đã phát hành, không trùng};
    \node[box, below=of filter] (derive) {6. Tạo biến mới, giữ các cột cần dùng};
    \node[io, below=of derive] (final) {Dữ liệu cho đề tài: \PreRowsFinal{} dòng $\times$
      \PreColsFinal{} cột};
    \foreach \a/\b in {raw/read, read/select, select/clean, clean/parse,
        parse/filter, filter/derive, derive/final}
      \draw[arrow] (\a) -- (\b);
  \end{tikzpicture}
  \caption{Sơ đồ khối quy trình tiền xử lý}
  \label{fig:preprocess}
\end{figure}

% ---------------------------------------------------------------------
\subsection{Đọc dữ liệu}

Mọi cột được đọc dưới dạng chuỗi, chưa chuyển đổi giá trị nào, để từng giá trị đều được
kiểm tra định dạng ở các bước sau. Mỗi dòng được gán mã \texttt{gpu\_id} là vị trí của
dòng trong tệp gốc để truy vết.

\begin{lstlisting}[language=Python, caption={Đọc dữ liệu thô dưới dạng chuỗi}]
raw = pd.read_csv(path, dtype="str", keep_default_na=False)
\end{lstlisting}

% ---------------------------------------------------------------------
\subsection{Chọn biến}

Trong \PreColsRaw{} biến gốc, nhóm chỉ giữ \PreColsSelected{} biến cần cho đề tài:
{\sloppy
\begin{itemize}
  \item biến mục tiêu \texttt{Max\_Power} và các biến giải thích \texttt{TMUs},
    \texttt{Core\_Speed}, \texttt{Memory}, \texttt{Memory\_Bandwidth}, \texttt{Process};
  \item các biến phân nhóm và định danh: \texttt{Name}, \texttt{Manufacturer},
    \texttt{Notebook\_GPU}, \texttt{Memory\_Type}, \texttt{Release\_Date},
    \texttt{Architecture};
  \item các biến chỉ dùng trong lúc xử lý: \texttt{Dedicated} (lọc GPU rời),
    \texttt{ROPs} (sửa lỗi ``(x$N$)'' ở bước xử lý định dạng), \texttt{Texture\_Rate} (kiểm tra lỗi và
    tính hiệu quả năng lượng), \texttt{Boost\_Clock}, \texttt{Memory\_Bus},
    \texttt{Memory\_Speed} (kiểm tra lỗi).
\end{itemize}
\par}
Các biến còn lại như cổng xuất hình, độ phân giải, phiên bản DirectX, giá bán hay bộ
nguồn khuyến nghị không liên quan tới câu hỏi nghiên cứu nên bị bỏ. Tỉ lệ khuyết không
được dùng làm tiêu chí chọn biến, vì chính biến mục tiêu \texttt{Max\_Power} cũng thiếu
khá nhiều.

% ---------------------------------------------------------------------
\subsection{Làm sạch chuỗi}

{\sloppy Các biến được chọn chứa nhiều giá trị không chuẩn: có các ký tự xuống dòng và khoảng
trắng thừa (ngày phát hành ghi dạng \verb|"\n01-Mar-2009 "|), ký hiệu \verb|"-"| thay
cho giá trị thiếu (\PreCoreSpeedDash{} giá trị của \texttt{Core\_Speed}, cùng
\PreCoreSpeedBlank{} giá trị chỉ chứa ký tự trắng) và \PreUnknownReleaseDate{} giá trị ghi
\verb|"Unknown Release Date"|. Nhóm xóa khoảng trắng hai đầu, rồi đổi các ký hiệu trên
thành giá trị thiếu chuẩn \texttt{NA}.\par}

\begin{lstlisting}[language=Python, caption={Xóa khoảng trắng và đánh dấu giá trị thiếu}]
text = selected.apply(lambda column: column.str.strip())
text = text.replace(["", "-"], pd.NA)
text["Release_Date"] = text["Release_Date"].replace(
    "Unknown Release Date", pd.NA)
\end{lstlisting}

\textbf{Nhận xét.} Nếu chỉ đếm dữ liệu trống như mặc định của \texttt{pandas}, tỉ lệ thiếu
của \texttt{Core\_Speed} và \texttt{Release\_Date} sẽ bị đánh giá thấp. Vì vậy mọi tỉ lệ
thiếu trong mục này đều được tính sau khi đã đổi các ký hiệu trên thành \texttt{NA}.

% ---------------------------------------------------------------------
\subsection{Xử lý định dạng}

Phần lớn biến số được ghi kèm đơn vị. Nhóm tách phần số và đổi sang kiểu số theo mẫu
định dạng riêng của từng biến (Bảng~\ref{tab:formats}); giá trị nào không khớp mẫu thì
chương trình dừng và báo lỗi. Cách xóa mọi ký tự không phải chữ số tuy đơn giản hơn
nhưng cho kết quả sai ở hai trường hợp: \verb|"800MB/sec"| còn lại 800 và bị hiểu là
800~GB/s thay vì 0,8~GB/s; \verb|"16 (x2)"| còn lại các chữ số ``162'' thay vì
$16 \times 2 = 32$.

\begin{table}[!htbp]
  \centering
  \caption{Một số định dạng gốc và kết quả chuyển đổi}
  \label{tab:formats}
  \small
  \begin{tabular}{>{\ttfamily}l >{\ttfamily}l l}
    \toprule
    \normalfont\textbf{Biến gốc} & \normalfont\textbf{Giá trị gốc} & \textbf{Sau chuyển đổi} \\
    \midrule
    Max\_Power        & 141 Watts       & 141 (W) \\
    Core\_Speed       & 738 MHz         & 738 (MHz) \\
    Memory\_Bandwidth & 64GB/sec        & 64 (GB/s) \\
    Memory\_Bandwidth & 800MB/sec       & 0,8 (GB/s) \\
    Process           & 55nm            & 55 (nm) \\
    ROPs              & 16 (x2)         & 32 \\
    Release\_Date     & 01-Mar-2009     & năm 2009 \\
    \bottomrule
  \end{tabular}
\end{table}

\textbf{Nhận xét.}
\begin{itemize}
  \item \PreBandwidthInMb{} dòng ghi băng thông theo MB/sec được đổi sang GB/s.
  \item Các biến định tính \texttt{Manufacturer} và \texttt{Memory\_Type} được chuyển sang
    kiểu phân loại (\texttt{category}, tương đương \texttt{factor} trong R) để dùng làm
    biến phân nhóm.
  \item Giá trị dạng ``$a$ (x$N$)'' (\PreRopsMultiplier{} dòng của \texttt{ROPs}) là phép
    nhân: Radeon RX~580 ghi \verb|"16 (x2)"| và có $16 \times 2 = 32$ ROPs.
  \item Ở đúng các dòng này, \texttt{TMUs} bị ghi chia cho $N$. Bằng chứng là tốc
    độ xử lý texture phải bằng TMUs nhân xung nhịp, nhưng \PreTextureMismatchBefore{} dòng
    không thỏa; sau khi nhân lại với $N$ thì còn \PreTextureMismatchAfter{} dòng. Ví dụ
    GTX 1080 Ti được ghi 56 TMUs, giá trị đúng là 224. Theo nguyên tắc làm sạch dữ liệu,
    giá trị sai được sửa khi tìm được giá trị đúng \cite{vandenbroeck2005}; ở đây giá
    trị đúng tính được từ chính các cột khác của dòng đó. \texttt{TMUs} (số đơn vị xử lý
    texture) phản ánh độ lớn của GPU và là một biến giải thích công suất trong mô hình
    hồi quy, nên nếu không sửa, các card đời 2016--2017 sẽ bị xem như card nhỏ hơn nhiều
    so với thực tế. \texttt{Memory\_Bus} cũng bị chia cho $N$ ở các dòng này và được sửa
    theo cùng cách.
\end{itemize}

\begin{lstlisting}[language=Python, caption={Sửa lỗi ``(x$N$)'': nhân lại TMUs và Memory\_Bus với $N$}]
multiplier = parse_multiplier(selected["ROPs"])  # "16 (x2)" -> 2, "64" -> 1
tmus = parse_number(selected["TMUs"]) * multiplier
memory_bus_bit = parse_number(selected["Memory_Bus"], "Bit") * multiplier
\end{lstlisting}

% ---------------------------------------------------------------------
\subsection{Lọc quan sát}
\label{sec:filter}

Đề tài nghiên cứu \emph{một GPU rời đã bán ra thị trường}, nên nhóm loại GPU tích hợp,
card chưa phát hành và các cấu hình nhiều GPU (ghép nhiều card SLI/CrossFire, card hai
chip như GTX~690 hay R9~295X2), vì công suất ghi cho các cấu hình này là của cả hệ.
Cấu hình nhiều GPU được nhận diện bằng các quy tắc trên tên sản phẩm: tên chứa ``SLI'',
``CrossFire'', ``X2''\ldots{} hoặc thuộc danh sách các card hai chip đã biết.

\begin{table}[!htbp]
  \centering
  \caption{Số dòng còn lại sau từng bước lọc}
  \label{tab:rows}
  \small
  \begin{tabular}{lr}
    \toprule
    \textbf{Bước} & \textbf{Số dòng} \\
    \midrule
    Dữ liệu gốc & \PreRowsRaw \\
    Chỉ giữ GPU rời & \PreRowsDedicated \\
    Loại cấu hình nhiều GPU & \PreRowsSingleGpu \\
    Loại card chưa phát hành & \PreRowsReleased \\
    Loại dòng trùng hoàn toàn & \PreRowsFinal \\
    \bottomrule
  \end{tabular}
\end{table}

\textbf{Nhận xét.} Bước đầu loại \PreIntegrated{} GPU tích hợp và \PreUnknownDedicated{}
dòng không rõ loại GPU. Dữ liệu còn \PreRowsFinal{} GPU rời, gồm \PreDesktop{} card desktop
và \PreLaptop{} card laptop.

% ---------------------------------------------------------------------
\subsection{Tạo biến mới}

Bước này tạo các biến mà nhiều phân tích phải dùng cùng một định nghĩa, rồi giữ lại
\PreColsFinal{} cột cho bộ dữ liệu cuối (Bảng~\ref{tab:columns}, kèm tỉ lệ khuyết của
từng cột):
\begin{itemize}
  \item \texttt{segment}: Desktop hoặc Laptop.
  \item \texttt{needs\_aux\_power}: công suất trên 75~W, mức tối đa một khe PCIe cấp được.
  \item \texttt{memory\_type\_group}: loại bộ nhớ, gộp các loại có dưới 30 card thành
    ``Other'' để mỗi nhóm đủ lớn khi so sánh.
  \item \texttt{texture\_rate\_per\_watt}: hiệu quả năng lượng (tốc độ xử lý texture trên
    mỗi watt).
\end{itemize}

\begin{table}[!htbp]
  \centering
  \caption{Các cột của bộ dữ liệu sau tiền xử lý và tỉ lệ khuyết}
  \label{tab:columns}
  \small
  \begin{tabular}{l >{\ttfamily}l r}
    \toprule
    \textbf{Vai trò} & \normalfont\textbf{Cột} & \textbf{Khuyết (\%)} \\
    \midrule
    Định danh & gpu\_id & \PreMissingPctGpuId \\
              & name & \PreMissingPctName \\
    \midrule
    Biến mục tiêu & max\_power\_w & \PreMissingPctMaxPowerW \\
    \midrule
    Biến tạo từ mục tiêu & needs\_aux\_power & \PreMissingPctNeedsAuxPower \\
                         & texture\_rate\_per\_watt & \PreMissingPctTextureRatePerWatt \\
    \midrule
    Biến giải thích & tmus & \PreMissingPctTmus \\
                    & core\_speed\_mhz & \PreMissingPctCoreSpeedMhz \\
                    & memory\_bandwidth\_gb\_s & \PreMissingPctMemoryBandwidthGbS \\
                    & memory\_mb & \PreMissingPctMemoryMb \\
                    & process\_nm & \PreMissingPctProcessNm \\
    \midrule
    Biến phân nhóm & manufacturer & \PreMissingPctManufacturer \\
                   & segment & \PreMissingPctSegment \\
                   & memory\_type & \PreMissingPctMemoryType \\
                   & memory\_type\_group & \PreMissingPctMemoryTypeGroup \\
    \midrule
    Thời gian & release\_year & \PreMissingPctReleaseYear \\
    Kiến trúc chip & architecture & \PreMissingPctArchitecture \\
    \bottomrule
  \end{tabular}
\end{table}

% ---------------------------------------------------------------------
\subsection{Xử lý dữ liệu khuyết}

Biến mục tiêu \texttt{max\_power\_w} thiếu ở \PrePowerMissing{} dòng
(\PrePowerMissingPct\%). Nhóm không điền giá trị thiếu, vì điền bằng trung bình
hay trung vị sẽ làm sai lệch phân phối; mỗi phân tích tự chọn các dòng có đủ biến cần
dùng.

Các biến giải thích cũng có giá trị khuyết, nhiều nhất là \texttt{core\_speed\_mhz}
(\PreMissingPctCoreSpeedMhz\%) và \texttt{tmus} (\PreMissingPctTmus\%)
(Bảng~\ref{tab:columns}), nên một phân tích dùng càng nhiều biến thì càng còn ít dòng.

\textbf{Nhận xét.} Card thiếu số liệu công suất là card đời cũ hơn (trung vị năm phát
hành \PreMissingMedianYear{} so với \PreObservedMedianYear{}, kiểm định Mann--Whitney
$p$ \PreMissingYearP) và có tỉ lệ AMD cao hơn (\PreMissingAmdPct\% so với
\PreObservedAmdPct\%, $p$ \PreMissingMakerP). Vì vậy các thống kê trên toàn mẫu đại diện
nhiều hơn cho card đời mới.

% ---------------------------------------------------------------------
\subsection{Các hàm chính và cấu trúc dữ liệu}

Quy trình được tóm tắt ở Hình~\ref{fig:preprocess}; Bảng~\ref{tab:functions} giải thích
các hàm chính. Dữ liệu được lưu dưới dạng bảng \texttt{pandas.DataFrame}, mỗi dòng là một GPU,
mỗi cột là một biến.

\begin{table}[!htbp]
  \centering
  \caption{Các hàm chính trong quy trình tiền xử lý}
  \label{tab:functions}
  \small
  \begin{tabularx}{\linewidth}{>{\ttfamily}l X}
    \toprule
    \normalfont\textbf{Hàm} & \textbf{Chức năng} \\
    \midrule
    read\_raw & Đọc tệp gốc, mọi cột dưới dạng chuỗi \\
    select\_columns & Giữ các biến cần cho đề tài \\
    clean\_text & Xóa khoảng trắng, đánh dấu giá trị thiếu \\
    parse\_columns & Tách đơn vị, đổi kiểu, sửa lỗi ``(x$N$)''; báo lỗi nếu gặp định dạng lạ \\
    check\_consistency & Kiểm tra giá trị bằng đẳng thức giữa các biến \\
    filter\_rows & Lọc quan sát, đếm số dòng còn lại sau mỗi bước \\
    derive\_variables & Tạo biến mới, giữ các cột cần dùng \\
    \bottomrule
  \end{tabularx}
\end{table}

\textbf{Kết quả.} Bộ dữ liệu sau tiền xử lý gồm \PreRowsFinal{} GPU rời và
\PreColsFinal{} cột. Trong các dòng mẫu ở Bảng~\ref{tab:sample}, mỗi biến số đã ở dạng
số với một đơn vị, giá trị thiếu ghi là \texttt{NA}, và GTX 1080 Ti có số TMUs đúng.

\begin{table}[!htbp]
  \centering
  \caption{Một vài dòng của bộ dữ liệu sau tiền xử lý}
  \label{tab:sample}
  \footnotesize
  \setlength{\tabcolsep}{4pt}
  \begin{tabular}{l r r r r r r r}
    \toprule
    \textbf{Tên} & \textbf{Năm} & \textbf{Công suất} & \textbf{TMUs} & \textbf{Xung nhịp}
      & \textbf{Băng thông} & \textbf{Bộ nhớ} & \textbf{Tiến trình} \\
    & & (W) & & (MHz) & (GB/s) & (MB) & (nm) \\
    \midrule
    \csname @@input\endcsname generated/preprocessing_sample_rows.tex
    \bottomrule
  \end{tabular}
\end{table}
